\begin{abstract}
\noindent
Small language models are increasingly trained by distilling from larger models, and
several such models are often available. The student must then decide which one to learn
from, on which domain, and when; the common default is the strongest throughout. This
thesis treats the decision as a curriculum learning problem over pairs of provider and
domain, termed provider-selective distillation, in which the next pair is chosen from the
progress the student actually makes. A preliminary study with text-only providers and a
price on every query, using a 1.5B-parameter student, five providers and three domains,
finds that providers of very different strength teach at a similar
rate per gradient step, so that price decided value and the strongest provider was the
worst single choice. A selector that learns where progress happens matches the best single
provider (test utility 0.500 against 0.494) without being told which one it is. Learning
was confined to one domain and unstable. The thesis then studies a refined
problem without a budget, in which the student reads the providers' token probabilities,
and proposes a training framework that combines a ground-truth reward with supervision
from the selected provider, with the aim of a student that passes its teachers.

\vspace{5mm}
Subject Descriptors:
\begin{list}{\quad}{}
    \item I.2.6 Learning
    \item I.2.7 Natural Language Processing
    \item G.3 Probability and Statistics
\end{list}

Keywords: \\
\tab On-policy distillation, curriculum learning, multi-teacher distillation,
multi-armed bandits

Implementation Software and Hardware:
\begin{list}{\quad}{}
    \item Python 3.12, PyTorch, vLLM, TRL, PEFT, Math-Verify, bubblewrap
    \item NVIDIA L40 GPUs (shared server); SoC Slurm cluster (A100, H100)
\end{list}
\end{abstract}
\pagebreak
